% !TeX root = ../main.tex


\chapter{Conclusion}
\label{ch:conclusion}

This dissertation asked whether multitask learning helps predict the next category and
the next region of a visit, and which design choices determine the answer. Three studies
addressed this question in sequence. The first developed a joint model that did not
consistently outperform its dedicated counterparts (Chapter~\ref{ch:cbic}). The second
kept the model architecture fixed, changed the input representation, and identified that
representation as the main bottleneck in the tested configuration
(Chapter~\ref{ch:courb}). Those two studies paired next-category prediction with a static
category classification task, and the pair changed in the final study.
The third study combined a check-in-level representation with a redesigned sharing
topology (Chapter~\ref{ch:mobiwac}). Under this design, the joint model outperforms the
dedicated single-task models on the category task at all six datasets and on the region
task at four of the six. At the other two, it remains statistically non-inferior within
a two-point margin (TOST).


\section{Contributions by chapter}
\label{sec:conc:contributions}

Chapter~\ref{ch:cbic} presents MTLnet, the first joint model developed in this research.
It combined static category classification with next-category prediction and shared its
lower layers across the two tasks. Both tasks read a place-level graph embedding. In
this configuration, the joint model cost more to train and did not consistently
outperform the dedicated single-task models on either task. The chapter contributes
this negative finding and an analysis of three possible causes: task dissimilarity and
the resulting risk of negative transfer, an input representation that may be
insufficient for both tasks, and an overly rigid sharing topology. The finding is
limited to a place-level input under hard sharing. The next two studies examine what
happens when these conditions change.

Chapter~\ref{ch:courb} kept the MTLnet architecture fixed but replaced its single
64-dimensional place embedding with separate spatial, temporal, and categorical
encoders. On the static task, category macro-F1 rose by 20.2 to 22.0 percentage points
across the three states tested. This increase exceeded those obtained from the
architectural variations evaluated in the first study.

This gain requires two qualifications: First, the static task classifies a place from
that place's own representation, so the input already determines the target. The gain
therefore provides no evidence about the sequential task. Second, the comparison is not
width-matched: the decomposed input has 192 dimensions, whereas the place embedding has
64. Chapter~\ref{ch:courb} therefore calls for an equal-dimension control to separate the
semantic contribution of the encoders from the effect of the additional width.

The sequential task provides the relevant diagnosis because its input does not
determine its target. The decomposed encoders are more effective in most of the tested
cases. Because the sharing topology remained fixed, this result identifies the input
representation as the main bottleneck in that configuration.

Chapter~\ref{ch:mobiwac} responds to this diagnosis with Check2HGI, a check-in-level
representation that assigns a distinct vector to each visit rather than reusing one
vector for every visit to the same place. The joint model exchanges information between
two task-specific streams through cross-attention, while a private spatial path serves
the region task. The evaluation used user-disjoint cross-validation, with twenty fitted
models per configuration: four seeds, each drawing its own five-fold user partition.
Paired tests used the four per-seed means.

Under this protocol, the joint model outperforms the dedicated single-task models on
the category task at all six datasets, by 5.3 to 9.4 macro-F1 points. On the region task,
it outperforms them at four datasets: Istanbul, Florida, California, and Texas. At
Alabama and Arizona, it remains statistically non-inferior within a two-point margin
(TOST). Non-inferiority was the registered primary analysis for the region task; the
four superiority findings are secondary results outside that analysis plan. Further analyses
clarify what these results do and do not establish. They isolate the contribution of the
input representation, test parameter count as an alternative explanation for the next-category
gains, and identify a possible relation between problem scale and the next-region gains.

Beginning with input representation, Section~\ref{sec:mobiwac:results-part1} provides the first
explanation through a controlled comparison between Check2HGI and HGI. With the dedicated
model, training configuration, windowing, and folds held fixed, Check2HGI produces higher
next-category macro-F1 than the place embedding HGI across the six evaluated datasets.
This result establishes the input representation as one condition that affects
performance in this setting. It also supports testing Check2HGI in other mobility
prediction architectures, although its benefit in those architectures has not yet been
evaluated.

The joint model's larger capacity offers a second possible explanation for the
next-category gains. Appendix~\ref{apx:parameter-count-control} compares it with a
capacity-matched dedicated model. In Alabama and California, the wider model
produces slightly lower next-category macro-F1 than the original dedicated model and
remains below the joint model. Thus, parameter count alone does not explain the joint
model's next-category advantage in these two datasets. The comparison instead points to
the joint architecture or its multitask training, but it does not separate their effects.

For next-region prediction, the results point to problem scale as another possible
condition. Across the five U.S. states, the joint model's advantage over the dedicated
model grows with the number of regions. It exceeds two Acc@10 points in Texas and
California, the states with the largest number of regions. This pattern suggests that
joint learning becomes more useful when the model must choose among more regions.
However, states with more regions also tend to have more check-ins. A controlled
experiment is therefore required to separate the effect of region count from the amount
of data and other differences between the states. Problem scale remains a possible
condition, not an established cause.



\section{The consolidated answer}
\label{sec:conc:answer}

The research question asks whether multitask learning helps next-category and
next-region prediction, and what the answer depends on. For the first part, the evidence
shows that it helps under the final configuration. The sequence that led to this
configuration also identifies the conditions behind that answer. Identifying these
conditions is the main finding of this dissertation.

The explanation begins with Chapter~\ref{ch:cbic}, which used a place-level embedding and
naive hard sharing. Its joint model did not consistently outperform the two dedicated
models, leaving three possible explanations: the tasks might be too different, the input
might be insufficient, or the sharing topology might be too restrictive.
Chapter~\ref{ch:courb} separated one of these explanations by
keeping the architecture fixed and changing the input. Performance then improved,
including on the sequential task in most states. Because the architecture did not
change, this comparison identifies the input representation as the main bottleneck in
that configuration. This diagnosis set the direction for the final study.

Chapter~\ref{ch:mobiwac} responded to this diagnosis in two steps. First, it introduced
Check2HGI, which replaced one vector per place with one vector per visit. Second, the
joint architecture replaced naive hard sharing with task-specific streams that exchange
information through cross-attention, while the region output retained a private spatial
path.

Together, these changes define the final configuration evaluated in this dissertation.
They also explain why the answer is conditional. The initial premise was incomplete, explained in
Section~\ref{sec:fund:mtl}, where multitask learning trains related tasks together in
the expectation that a shared representation will improve generalization, was disproven by the three
studies showing that task relatedness and joint training are not sufficient by themselves.
The controlled comparisons identify the input representation as one condition, while suggests that the architecture
and dataset scale may matter, although their
individual effects were not fully isolated. The answer therefore depends on the design and
evaluation conditions, not on joint training alone.

These findings support a conditional rather than universal conclusion.
The consolidated answer is therefore not that multitask learning always helps.
It is that multitask learning helps next-category and next-region prediction under
the final design and evaluation protocol developed in this dissertation.

\section{Limitations}
\label{sec:conc:limitations}

Six limitations bound the scope of these conclusions.

\begin{enumerate}
    \item \textbf{Data vintage.} The five state datasets come from
    Gowalla~\cite{cho2011gowalla}, and the extraction used here spans January 2009
    to August 2011 across the five states. The Istanbul
    dataset~\cite{wongso2025massivesteps} is not appreciably more recent for most of its
    volume: its check-ins fall in two separate periods, 2012 to 2013 and 2017 to 2018,
    with none in between, and roughly seven in ten belong to the earlier period.
    Mobility patterns, place inventories,
    and check-in behavior have changed since.
    \item \textbf{Taxonomy coarseness.} Next-category prediction uses seven top-level
    classes. A finer taxonomy may change the effect of joint training.
    \item \textbf{Transductive representation.} Check2HGI is trained on the check-in
    graph of each dataset, so it cannot represent unseen places or users without
    retraining. The place-level representation it builds
    on~\cite{huang2023hgi} is trained the same way, and methods that learn aggregators
    instead of per-node vectors~\cite{hamilton2017graphsage} are the alternative this
    constraint points toward.
    \item \textbf{No next-place task.} The experiments do not predict the exact next
    POI, and their conclusions apply only to next category and next region.
    \item \textbf{Geographic coverage.} Outside the United States, the evidence rests on
    a single city, Istanbul.
    \item \textbf{The task-pair confound.} The task pair changed together with the
    representation and the sharing topology. Chapters~\ref{ch:cbic}
    and~\ref{ch:courb} paired static category classification with next-category
    prediction, while Chapter~\ref{ch:mobiwac} pairs two sequential targets,
    next-category and next-region prediction. No single controlled ablation
    separates the representation-and-topology change from the task-pair change in
    the final result. Chapter~\ref{ch:courb} is the fixed-pair control for the
    diagnosis. The additional parameter-count experiment in
    Appendix~\ref{apx:parameter-count-control} tests only whether a larger dedicated
    category model can close the gap. It covers two of the six datasets, one width
    point per dataset, width scaling rather than depth, and a smaller hyperparameter
    search than the original dedicated model. The greater similarity of
    the final pair may still contribute to the size of the improvement. The
    ablation that would separate the two changes,
    running static category classification under the check-in-level representation, is
    not clean under that representation: the category of the visited place is an input
    feature of a check-in node, so that target would be partly readable from its own
    input. This follows from the design of the representation and was not measured, so
    the confound is bounded by the fixed-pair control rather than removed.
\end{enumerate}


\section{Future work}
\label{sec:conc:future}

Future work follows directly from these limitations. Newer and denser traces would
test the conclusions beyond the Gowalla vintage (limitation~1), and finer-grained
taxonomies would test them beyond the seven-class division (limitation~2). Because the
check-in-level representation is trained on a fixed graph, an inductive variant of it,
one that embeds unseen places and users without retraining, would remove the
transductive constraint and support deployment in growing cities. Two further changes to
the same representation would test how much of its behavior follows from how it is
assembled rather than from the hierarchy: its place-vector table is initialized from a
pretrained table and then held near it during training, so an ablation that varies that
coupling would separate the contribution of the initialization, and a hypergraph
formulation, in which one edge joins the several visits of a session rather than only
consecutive pairs, would test whether the pairwise graph is the right structure for a
visit history (limitation~3).

Adding the exact next place as a third target would extend the joint model beyond the
two properties predicted in this work. The cascade architectures reviewed in
Chapter~\ref{ch:mobiwac} suggest that category and region predictions can provide
additional structure for this task, and reusing the existing check-in-level
representation would require a change to the input pipeline and one additional output
rather than a new representation, which is also the route by which this representation
would serve contexts other than the two studied here (limitation~4).
Further cities outside the United States would widen the geographic base
(limitation~5).
Isolating the effect of changing the task pair needs a static target
that the check-in-level representation does not already carry as an input feature, since
running the earlier pair directly under this representation is not a clean comparison
(limitation~6).



\section{Final remarks}
\label{sec:conc:final}

The practical output of this dissertation is a single model that predicts two
properties of the next visit in one forward pass: what kind of place the user will
visit and in which part of the city. The methodological contribution is the sequence
of evidence that led to this model: a published negative result, the diagnosis that
identified the input representation as the main bottleneck, and a solution designed
from that diagnosis. The negative result was not an obstacle to the contribution. It
was its first half.
